\section{The genuine homogeneous action on Rees cohomology}

\begin{lemma}[Homogeneous scalars preserve the original degree shift]
\label{lem:fc-rees-cech-homogeneous-action}
\leanok
\lean{EnvelopeProperty.FcReesCechHomogeneousAction.chartCoefficient_monomial, EnvelopeProperty.FcReesCechHomogeneousAction.termPowerEquiv_symm_lof, EnvelopeProperty.FcReesCechHomogeneousAction.homogeneousScalar, EnvelopeProperty.FcReesCechHomogeneousAction.term_homogeneous, EnvelopeProperty.FcReesCechHomogeneousAction.gradedHomogeneous, EnvelopeProperty.FcReesCechHomogeneousAction.degreeInclusion_homogeneous, EnvelopeProperty.FcReesCechHomogeneousAction.homologyInclusion_homogeneous, EnvelopeProperty.FcReesCechHomogeneousAction.ordered_homogeneous, EnvelopeProperty.FcReesCechHomogeneousAction.originalHomologyIso, EnvelopeProperty.FcReesCechHomogeneousAction.originalHomologyIso_homogeneous, EnvelopeProperty.FcReesCechHomogeneousAction.reesModule_homogeneous}
\uses{lem:fc-rees-cech-degree-inclusions,lem:fc-rees-cech-homogeneous-multiplication,lem:fc-rees-cech-homology-module}
Let $X\to\operatorname{Spec}R$ be separated and let $(U_i)$ be a finite
linearly ordered family of affine opens.  For $a\in I^e$, multiplication by
the genuine Rees monomial $at^e$ on the structure-module Čech complex of
the schematic Rees family carries the original degree-$d$ summand into
the original degree-$(e+d)$ summand.  Under the canonical cohomology
decomposition its action on that summand is the cohomology map induced
by multiplication by $a$ on the original ideal-power sections.
\end{lemma}
\begin{proof}
The coefficient decomposition sends the polynomial $xt^d$ to the
degree-$d$ inclusion of $x$, on every original intersection.  Conversely,
polynomial coefficients show that the inverse image of this inclusion
is exactly $xt^d$.  The original coefficient map sends $at^e$ to
$(a|_V)t^e$ on an affine intersection $V$, where $a|_V$ is its original
structural coefficient.  Thus its product with $xt^d$ is
$(a|_V)x t^{e+d}$, at every intersection.  This proves equality of the
entire term maps.  Homogeneous multiplication commutes with the
original restrictions, so these maps give the equality of cochain maps.
Applying cohomology preserves their equality and compositions.  The
cohomology decomposition preserves each original summand inclusion;
linearity of the structure-module comparison and of the homology
functor therefore identifies the genuine scalar action with the stated
degree-shift map on every cohomology class.
\end{proof}
